% !TEX root = ../algorithms.tex

	\section{Ранжирование списка во внешней памяти. Линейное программирование}
	
	\subsection{Задача}
	Дан список, хранящийся в массиве $A[0..N-1]$ в произвольном порядке.
	У каждого узла $x$ есть поля
	\[
	x.\texttt{next},\qquad x.\texttt{prev},\qquad x.w=1.
	\]
	Также в реализации удобно хранить поле
	\[
	x.\texttt{id},
	\]
	причём изначально
	\[
	A[i].\texttt{id}=i.
	\]

	
	Для узла \(x\) хотим получить его \emph{ранг}
	\[
	\mathrm{rank}(x)=\text{число переходов по } \mathrm{next} \text{ от } x \text{ до конца списка}.
	\]

	
	\subsection{Идея сжатия списка}
	
	Пусть \(S\) --- множество удаляемых узлов, такое что никакие два соседних по списку узла не удаляются одновременно.
	\[
	\forall x\in S\quad x.\mathrm{prev} \notin S \text{ и } |S| = \Theta(N).
	\]
	
	\medskip
	
	Тогда все узлы из \(S\) можно удалить \emph{независимо друг от друга}:
	 \(\forall x\in S\) выполняем
	\[
	x.\mathrm{prev}.\mathrm{next} = x.\mathrm{next},\qquad
	x.\mathrm{prev}.\mathrm{w} += x.\mathrm{w}.
	\]
	
	После удаления получаем более короткий список. Затем задача рекурсивно решается на нём. После возврата из рекурсии удалённые узлы нужно восстановить и дописать для них ответы.
	
	\subsection{Как выбрать множество \(S\)}
	
	Каждой вершине независимо присваиваем случайный бит
	\[
	x.r\in\{0,1\},
	\qquad
	\Pr[x.r=0]=\Pr[x.r=1]=\frac12.
	\]
	
	После этого выбираем
	\[
	x\in S
	\iff
	x.r=1 \ \text{и}\ x.\texttt{prev}.r=0.
	\]
	
	То есть в $S$ попадают вершины, которые являются ``правым концом'' шаблона $01$ вдоль списка.
	
	\subsubsection*{Почему это корректно}
	
	Если $x\in S$, то по определению
	\[
	x.\texttt{prev}.r=0,\qquad x.r=1.
	\]
	Значит, $x.\texttt{prev}$ не может одновременно лежать в $S$, потому что для этого у него должен был бы быть бит $1$.
	
	Следовательно,
	\[
	\forall x\in S\colon x.\texttt{prev}\notin S.
	\]
	
	\subsubsection*{Размер множества $S$}
	
	Для внутренней вершины вероятность попасть в $S$ равна
	\[
	\Pr[x\in S]=\Pr[x.\texttt{prev}.r=0 \land x.r=1]=\frac14.
	\]
	
	Поэтому
	\[
	\mathbb{E}|S| \approx \frac{N}{4},
	\]
	а после одного шага остаётся примерно
	\[
	N-|S| \approx \frac{3N}{4}
	\]
	вершин.
	
	\subsection{Рекурсивная схема алгоритма}
	
	\begin{algorithm}[H]
		\caption{\textsc{ListRank}$(A)$}
		\begin{algorithmic}[1]
			\Require Массив вершин списка $A$
			\Ensure Обновлённые значения $w$ для вершин списка
			
			\If{$|A| \le 1$}
			\State \Return $A$
			\EndIf
			
			\State Для каждой вершины $x \in A$ независимо выбрать случайный бит $x.r \in \{0,1\}$
			\State $S = \{\, x \in A \mid x.\mathrm{prev} \neq \textsc{null},\ x.\mathrm{prev}.r = 0,\ x.r = 1 \,\}$
			
			\Comment{Сжимаем список, удаляя вершины из $S$}
			\ForAll{$x \in S$}
			\State $p = x.\mathrm{prev}$
			\State $p.w = p.w + x.w$
			\State $p.\mathrm{next} = x.\mathrm{next}$
			\EndFor
			
			\State $A' = A \setminus S$
			\State \Call{ListRank}{$A'$}
			
			\Comment{Восстанавливаем удалённые вершины в обратном порядке}
			\ForAll{$x \in S$ в порядке, обратном удалению}
			\State $p = x.\mathrm{prev}$
			\State $p.\mathrm{next} = x$
			\State $p.w = p.w - x.w$
			\State $x.w = x.w + p.w$
			\EndFor
			
			\State \Return $A$
		\end{algorithmic}
		\end{algorithm}
	
	\subsection{Реализация через сортировки во внешней памяти}
	
	Главная цель --- заменить случайные обращения к указателям на несколько сортировок и линейных проходов.
	
	\subsubsection*{Храним идентификаторы}
	
	У каждого узла есть поле
	\[
	x.\texttt{id},
	\qquad
	A[i].\texttt{id}=i.
	\]
	
	Это позволяет после перестановок и сортировок восстановить ``исходную'' вершину.
	
	\subsubsection*{Сортируем по \texttt{next}}
	
	Пусть массив $A$ отсортирован по \texttt{id}. Тогда строим
	\[
	\widetilde A = \Sort(\{x\in A\}\text{ по }x.\texttt{next}).
	\]
	
	После такой сортировки запись $\widetilde A[i]$ стоит напротив той вершины, в которую она указывает.
	Поэтому при одновременном просмотре массивов $A$ и $\widetilde A$ можно сопоставить вершину и её предшественника.
	
	\subsubsection*{Находим удаляемые вершины и сразу обновляем предшественников}
	
	При линейном проходе по $A$ и $\widetilde A$:
	
	\[
	\text{если } \widetilde A[i].r=0 \text{ и } A[i].r=1,
	\text{ то удаляем } A[i].
	\]
	
	При этом выполняются обновления:
	\[
	\widetilde A[i].w \mathrel{+}= A[i].w,
	\qquad
	\widetilde A[i].\texttt{next}=A[i].\texttt{next},
	\]
	а сама удалённая вершина сохраняется в массив $S$.
	
	\subsubsection*{Убираем ``дубликаты'' после перенаправления указателей}
	
	После удаления вершины её предшественник начинает указывать в того же преемника, в которого раньше указывала удалённая вершина.
	Поэтому после повторной сортировки по \texttt{next}
	\[
	\widetilde A = \Sort(\widetilde A\text{ по } \texttt{next})
	\]
	появляются соседние элементы с одинаковым полем \texttt{next}.
	
	Из каждой такой пары надо оставить только живую вершину, а удалённую выбросить.
	Это делается линейным проходом по отсортированному массиву:
	если у двух соседних элементов одинаковый \texttt{next}, то удаляется тот, у которого $r=1$.
	
	После этого выжившие вершины упаковываются в плотный массив длины $N-|S|$.
	
	\subsubsection*{Перенумерация и рекурсия}
	
	После упаковки нужно переопределить индексы \texttt{next} под новый массив,
	затем отсортировать вершины по \texttt{id} и рекурсивно решить задачу уже для массива размера $N-|S|$:
	\[
	\widehat A = \Sort(\{x\in A[0..N-|S|-1]\}\text{ по }x.\texttt{id}).
	\]
	
	\subsubsection*{Восстановление}
	
	Предполагается, что:
	\begin{itemize}
		\item $S$ хранится в порядке возрастания \texttt{id};
		\item $\widehat A$ --- результат рекурсивного решения на сжатом массиве длины $N-|S|$.
	\end{itemize}
	
	Сначала сортируем удалённые и оставшиеся вершины по полю \texttt{next}:
	\[
	\widetilde S = \Sort(\{x\in S\}\text{ по }x.\texttt{next}),
	\qquad
	\widetilde A = \Sort(\{x\in \widehat A[0..N-|S|-1]\}\text{ по }x.\texttt{next}).
	\]
	
	После этого выполняем линейный проход:
	\begin{algorithm}[H]
		\caption{\textsc{Restore}$(\widehat A, S)$}
		\begin{algorithmic}[1]
			\State $\widetilde S = \Sort(S \text{ по } x.\texttt{next})$
			\State $\widetilde A = \Sort(\widehat A[0..N-|S|-1] \text{ по } x.\texttt{next})$
			
			\State $i = 0$; $j = 0$; $k = 0$; $m = 0$
			\While{$j < N-|S|$}
			\While{$k < |S| \ \textbf{and}\ i = S[k].\texttt{id}$}
			\State $i = i + 1$
			\State $k = k + 1$
			\EndWhile
			
			\State $\widetilde A[j].\texttt{next} = i$
			
			\If{$m < |S| \ \textbf{and}\ \widetilde A[j].\texttt{next} = \widetilde S[m].\texttt{next}$}
			\State $\widetilde A[j].\texttt{next} = \widetilde S[m].\texttt{id}$
			\State $\widetilde A[j].w = \widetilde A[j].w - \widetilde S[m].w$
			\State $\widetilde S[m].w = \widetilde S[m].w + \widetilde A[j].w$
			\State $m = m + 1$
			\EndIf
			
			\State $i = i + 1$
			\State $j = j + 1$
			\EndWhile
			
			\State $A = \Sort(x \in (\widetilde S \cup \widetilde A[0..N-|S|-1]) \text{ по } \texttt{id})$
			\State \Return $A$
		\end{algorithmic}
	\end{algorithm}
	
	\subsection{Оценка сложности}
	
	Обозначим через $\Sort(N)$ стоимость сортировки $N$ записей во внешней памяти.
	
	Один уровень рекурсии состоит из:
	\begin{itemize}
		\item константного числа сортировок по полям \texttt{id} и \texttt{next};
		\item константного числа линейных проходов по массивам.
	\end{itemize}
	
	Поэтому стоимость одного уровня равна
	\[
	O(\Sort(N)).
	\]
	
	\medskip
	
	При случайном выборе множества $S$ каждая внутренняя вершина попадает в $S$ с вероятностью $1/4$, поэтому на одном шаге удаляется в среднем постоянная доля вершин. Значит, размер задачи после сжатия в среднем уменьшается в постоянное число раз.
	
	Если обозначить через $T(N)$ время работы алгоритма на списке из $N$ вершин, то естественная рекурсия для математического ожидания имеет вид
	\[
	\mathbb{E}[T(N)] = \mathbb{E}[T(\alpha N)] + O(\Sort(N)).
	\qquad \alpha < 1,
	\]

	
	\subsubsection*{Доказательство утверждения}
	
	Разворачивая рекурсию, получаем
	\[
	\mathbb{E}[T(N)]
	=
	O(\Sort(N))
	+
	O(\Sort(\alpha N))
	+
	O(\Sort(\alpha^2 N))
	+\dots
	\]
	
	Размеры задач образуют геометрическую прогрессию:
	\[
	N,\ \alpha N,\ \alpha^2 N,\ \dots
	\]
	поэтому глубина рекурсии равна $O(\log N)$.
	
	Теперь подставим стандартную оценку для сортировки во внешней памяти:
	\[
	\Sort(N)=\Theta\!\left(
	\frac{N}{B}\log_{M/B}\frac{N}{B}
	\right).
	\]
	
	Тогда
	\[
	\sum_{k\ge 0}\Sort(\alpha^k N)
	=
	O\!\left(
	\sum_{k\ge 0}
	\frac{\alpha^k N}{B}
	\log_{M/B}\frac{\alpha^k N}{B}
	\right).
	\]
	
	Так как $\alpha^k N \le N$, имеем
	\[
	\log_{M/B}\frac{\alpha^k N}{B}
	\le
	\log_{M/B}\frac{N}{B},
	\]
	и потому
	\[
	\sum_{k\ge 0}\Sort(\alpha^k N)
	\le
	O\!\left(
	\log_{M/B}\frac{N}{B}
	\sum_{k\ge 0}\frac{\alpha^k N}{B}
	\right).
	\]
	
	Сумма геометрической прогрессии
	\[
	\sum_{k\ge 0}\alpha^k
	\]
	ограничена константой, следовательно,
	\[
	\sum_{k\ge 0}\Sort(\alpha^k N)
	=
	O\!\left(
	\frac{N}{B}\log_{M/B}\frac{N}{B}
	\right)
	=
	O(\Sort(N)).
	\]
	
	Итак,
	\[
	\mathbb{E}[T(N)] = O(\Sort(N)).
	\]
	
	\subsection*{Итог}
	
	Таким образом, ожидаемая сложность алгоритма равна
	\[
	\mathbb{E}[T(N)] = O(\Sort(N))
	=
	O\!\left(
	\frac{N}{B}\log_{M/B}\frac{N}{B}
	\right).
	\]
	
	\section*{Линейное программирование}
	
	\subsection{Общая постановка задачи ЛП}
	
	\textbf{Линейное программирование (ЛП)} --- это задача оптимизации линейной функции при линейных ограничениях.
	
	В общем виде задача может быть записана так:
	\[
	\max_{x \in \mathbb{R}^n} c^T x
	\qquad \text{при} \qquad
	Ax \le b,
	\]
	где
	\[
	A \in \mathbb{R}^{m \times n}, \qquad b \in \mathbb{R}^m, \qquad c \in \mathbb{R}^n.
	\]
	
	Неравенство \(Ax \le b\) понимается \textbf{покомпонентно}, то есть
	\[
	(Ax)_i \le b_i, \qquad i=1,\dots,m.
	\]
	
	То же самое в координатной записи:
	\[
	\max \; c_1x_1+\dots+c_nx_n
	\]
	при условиях
	\[
	\begin{cases}
		a_{11}x_1+\dots+a_{1n}x_n \le b_1,\\
		a_{21}x_1+\dots+a_{2n}x_n \le b_2,\\
		\hspace{4em}\vdots\\
		a_{m1}x_1+\dots+a_{mn}x_n \le b_m.
	\end{cases}
	\]
	
 	Две наиболее употребительные записи:
	
	\[
	\max \; c^T x \quad \text{при } Ax \le b,\; x \ge 0,
	\]
	и
	\[
	\min \; c^T x \quad \text{при } Ax = b,\; x \ge 0.
	\]
	
	Обе формы эквивалентны общей задаче после несложных преобразований.
	
	\subsection{Эквивалентные преобразования задач ЛП}
	
	Многие задачи линейного программирования можно привести к одной удобной форме.
	
	\paragraph{1) Минимизацию можно заменить максимизацией}
	\[
	\min \; c^T x
	\quad \Longleftrightarrow \quad
	\max \; (-c)^T x.
	\]
	
	\paragraph{2) Ограничение вида \texorpdfstring{$\ge$}{>=} можно заменить ограничением вида \texorpdfstring{$\le$}{<=}}
	Если
	\[
	a_k^T x \ge b_k,
	\]
	то, умножив обе части на \(-1\), получаем
	\[
	-a_k^T x \le -b_k.
	\]
	
	\paragraph{3) Равенство можно заменить двумя неравенствами}
	Если
	\[
	a_k^T x = b_k,
	\]
	то это эквивалентно системе
	\[
	\begin{cases}
		a_k^T x \le b_k,\\
		-a_k^T x \le -b_k.
	\end{cases}
	\]
	
	\paragraph{4) Свободную переменную можно заменить разностью двух неотрицательных}
	Если переменная \(x_j\) может быть любого знака, то вводят
	\[
	x_j^+ \ge 0, \qquad x_j^- \ge 0,
	\]
	и полагают
	\[
	x_j = x_j^+ - x_j^-.
	\]
	
	Это стандартный способ перейти к задаче с неотрицательными переменными.
	
	\paragraph{5) Неравенство \texorpdfstring{$\le$}{<=} можно превратить в равенство}
	Если имеется ограничение
	\[
	a_k^T x \le b_k,
	\]
	то вводится \textbf{добавочная переменная} (slack variable)
	\[
	s_k \ge 0,
	\]
	такая что
	\[
	a_k^T x + s_k = b_k.
	\]
	
	Смысл переменной \(s_k\): это \emph{остаток ресурса} или \emph{запас} до границы ограничения.
	
	\subsection{Приведение к канонической форме}
	
	Одна из самых удобных форм:
	\[
	\min \; c^T x
	\qquad \text{при} \qquad
	Ax = b, \qquad x \ge 0.
	\]
	
	Чтобы привести к ней произвольную линейную задачу, обычно делают так:
	
	\begin{enumerate}
		\item при необходимости заменяют максимум на минимум или наоборот;
		\item ограничения вида \(\ge\) заменяют на \(\le\), умножая на \(-1\);
		\item неравенства заменяют равенствами с помощью добавочных переменных;
		\item каждую свободную переменную заменяют разностью двух неотрицательных.
	\end{enumerate}
	
	\subsection{Геометрическая интерпретация}
	
	Пусть строки матрицы \(A\) равны
	\[
	a_1^T,\dots,a_m^T.
	\]
	Тогда каждое ограничение
	\[
	a_k^T x \le b_k
	\]
	задаёт \textbf{полупространство}.
	
	Следовательно, допустимое множество
	\[
	P=\{x \in \mathbb{R}^n \mid Ax \le b\}
	\]
	есть пересечение полупространств, то есть \textbf{выпуклый многогранник}.
	
	Целевая функция
	\[
	c^T x
	\]
	задаёт семейство гиперплоскостей
	\[
	c^T x = d.
	\]
	При изменении \(d\) эта гиперплоскость сдвигается параллельно самой себе. Задача ЛП состоит в том, чтобы найти такую гиперплоскость, которая:
	\begin{enumerate}
		\item ещё пересекает допустимое множество;
		\item при этом даёт наибольшее (или наименьшее) значение целевой функции.
	\end{enumerate}
	
	\begin{Example}{План производства}{}
	
	Предприятие выпускает два вида продукции:
	\[
	x_1 = \text{число столов}, \qquad x_2 = \text{число стульев}.
	\]
	
	Прибыль:
	\[
	40 \text{ у.е. за стол}, \qquad 30 \text{ у.е. за стул}.
	\]
	
	Есть два ресурса:
	\begin{itemize}
		\item древесина: не более \(18\) единиц;
		\item рабочее время: не более \(10\) часов.
	\end{itemize}
	
	На производство одной единицы продукции требуется:
	\[
	\begin{array}{c|cc}
		& \text{стол } (x_1) & \text{стул } (x_2)\\
		\hline
		\text{древесина} & 3 & 2\\
		\text{время} & 2 & 1
	\end{array}
	\]
	
	Тогда задача ЛП имеет вид:
	\[
	\max \; 40x_1 + 30x_2
	\]
	при ограничениях
	\[
	\begin{cases}
		3x_1 + 2x_2 \le 18,\\
		2x_1 + x_2 \le 10,\\
		x_1 \ge 0,\; x_2 \ge 0.
	\end{cases}
	\]
	
	\textbf{Смысл ограничений:}
	\begin{itemize}
		\item \(3x_1+2x_2 \le 18\) --- нельзя израсходовать больше древесины, чем есть;
		\item \(2x_1+x_2 \le 10\) --- нельзя превысить имеющийся запас рабочего времени;
		\item \(x_1,x_2 \ge 0\) --- отрицательное число изделий производить нельзя.
	\end{itemize}
	\end{Example}
	
	\subsection{Приближённое решение системы \texorpdfstring{$Ax \approx b$}{Ax ≈ b}}
	
	Иногда систему
	\[
	Ax=b
	\]
	нельзя решить точно: например, она переопределена или содержит ошибки измерений. Тогда ищут \(x\), для которого невязка
	\[
	Ax-b
	\]
	как можно меньше.
	
	Рассмотрим две постановки:
	\[
	\min_x \|Ax-b\|_2
	\qquad \text{и} \qquad
	\min_x \|Ax-b\|_1.
	\]
	
	\subsubsection*{Норма \texorpdfstring{$\ell_2$}{l2}: метод наименьших квадратов}
	
	Задача
	\[
	\min_x \|Ax-b\|_2
	\]
	имеет тот же набор решений, что и задача
	\[
	\min_x \|Ax-b\|_2^2,
	\]
	потому что квадрат --- монотонная функция на \([0,\infty)\).
	
	Именно для \textbf{квадрата нормы} выводятся \textbf{нормальные уравнения}:
	\[
	A^T A x = A^T b.
	\]
	
	\textbf{Почему получаются нормальные уравнения?}
	
	Рассмотрим функцию
	\[
	f(x)=\|Ax-b\|_2^2=(Ax-b)^T(Ax-b).
	\]
	
	Необходимое условие минимума: все частные производные по переменным
	\(x_1,\dots,x_n\) равны нулю.
	
	Вычисление этих производных приводит к равенству
	\[
	2A^T(Ax-b)=0.
	\]
	
	Следовательно,
	\[
	A^T(Ax-b)=0,
	\]
	то есть
	\[
	A^TAx=A^Tb.
	\]
	
	Это и есть нормальные уравнения.
	
	\medskip
	
	\textbf{Замечание.} Если столбцы матрицы \(A\) линейно независимы, то матрица \(A^T A\) невырождена, и решение единственно.
	
	\paragraph{Оценка сложности} для плотной матрицы \(A \in \mathbb{R}^{m\times n}\):
	\begin{itemize}
		\item вычисление \(A^T A\): \(O(mn^2)\);
		\item вычисление \(A^T b\): \(O(mn)\);
		\item решение системы размера \(n\times n\): \(O(n^3)\);
		\item память: \(O(n^2)\) для хранения \(A^T A\).
	\end{itemize}
	
	\subsubsection*{Норма \texorpdfstring{$\ell_1$}{l1}: сведение к линейному программированию}
	
	Задача
	\[
	\min_x \|Ax-b\|_1
	\]
	может быть сведена к задаче линейного программирования.
	
	Напомним, что
	\[
	\|Ax-b\|_1=\sum_{i=1}^m |(Ax-b)_i|.
	\]
	
	Чтобы убрать модуль, вводят новые переменные
	\[
	t_1,\dots,t_m \ge 0,
	\]
	где каждая \(t_i\) должна быть не меньше абсолютной величины соответствующей невязки:
	\[
	|(Ax-b)_i| \le t_i.
	\]
	
	Это эквивалентно двум линейным ограничениям:
	\[
	-(t_i) \le (Ax-b)_i \le t_i.
	\]
	
	Значит, задача переписывается так:
	\[
	\min_{x,t} \sum_{i=1}^m t_i
	\]
	при условиях
	\[
	Ax-b \le t,
	\qquad
	-(Ax-b) \le t,
	\qquad
	t \ge 0.
	\]
	
	Эквивалентно:
	\[
	\min_{x,t} \mathbf{1}^T t
	\]
	при
	\[
	\begin{cases}
		Ax - t \le b,\\
		-Ax - t \le -b,\\
		t \ge 0.
	\end{cases}
	\]
	
	\paragraph{Почему это работает}
	На оптимуме каждая переменная \(t_i\) станет минимально возможной, то есть
	\[
	t_i = |(Ax-b)_i|.
	\]
	Поэтому целевая функция
	\[
	\sum_{i=1}^m t_i
	\]
	ровно совпадает с \(\|Ax-b\|_1\).
	
	\paragraph{Размер новой задачи}
	\begin{itemize}
		\item новых переменных: \(m\);
		\item новых линейных ограничений: \(2m\);
		\item стоимость построения для плотного случая: \(O(mn)\).
	\end{itemize}
	
	\begin{Example}{ для \texorpdfstring{$\ell_1$}{l1}}{}
	Пусть
	\[
	A=
	\begin{pmatrix}
		1\\
		2
	\end{pmatrix},
	\qquad
	b=
	\begin{pmatrix}
		1\\
		4
	\end{pmatrix}.
	\]
	Тогда
	\[
	Ax-b=
	\begin{pmatrix}
		x-1\\
		2x-4
	\end{pmatrix},
	\]
	и задача
	\[
	\min_x \|Ax-b\|_1
	\]
	имеет вид
	\[
	\min_x \bigl(|x-1|+|2x-4|\bigr).
	\]
	
	Введём \(t_1,t_2 \ge 0\). Тогда получаем ЛП:
	\[
	\min_{x,t_1,t_2} t_1+t_2
	\]
	при условиях
	\[
	\begin{cases}
		x-1 \le t_1,\\
		-(x-1) \le t_1,\\
		2x-4 \le t_2,\\
		-(2x-4) \le t_2,\\
		t_1 \ge 0,\; t_2 \ge 0.
	\end{cases}
	\]
	
	Это уже обычная задача линейного программирования без модулей.
	\end{Example}
	
	\begin{Proposition}{}{}
		Как решить \texorpdfstring{$\min c^T x$}{min cTx} при \texorpdfstring{$Ax=b$}{Ax=b}, если переменные свободны по знаку	
	\end{Proposition}
	
	\subsubsection*{Идея решения}
	
	Если уже есть алгоритм, который умеет решать задачи вида
	\[
	\min \; \tilde c^T z
	\qquad \text{при} \qquad
	\tilde A z=\tilde b,\qquad z \ge 0,
	\]
	то нужно только избавиться от свободных переменных.
	
	Для каждого \(j=1,\dots,n\) вводим
	\[
	x_j=x_j^+-x_j^-,
	\qquad
	x_j^+ \ge 0,\quad x_j^- \ge 0.
	\]
	
	Тогда
	\[
	x=x^+-x^-,
	\]
	и исходная задача эквивалентна
	\[
	\min \; c^T(x^+-x^-)
	\]
	при
	\[
	A(x^+-x^-)=b,
	\qquad
	x^+ \ge 0,\quad x^- \ge 0.
	\]
	
	В матричном виде:
	\[
	\min
	\begin{pmatrix}
		c\\
		-c
	\end{pmatrix}^T
	\begin{pmatrix}
		x^+\\
		x^-
	\end{pmatrix}
	\]
	при
	\[
	\begin{pmatrix}
		A & -A
	\end{pmatrix}
	\begin{pmatrix}
		x^+\\
		x^-
	\end{pmatrix}
	=b,
	\qquad
	\begin{pmatrix}
		x^+\\
		x^-
	\end{pmatrix}
	\ge 0.
	\]
	
	Это уже задача в канонической форме.
	
	\subsubsection*{Размер новой задачи}
	\begin{itemize}
		\item было \(n\) переменных, стало \(2n\);
		\item число равенств осталось \(m\).
	\end{itemize}
	
	Построение новой задачи требует \(O(mn)\) операций для плотной матрицы.
	
	